\section{Literature positioning (minimal)}
\label{sec:lit}

The present work adopts the notion of a phase boundary from statistical physics
strictly as a \emph{structural analogy} for threshold-induced partitions of
admissible state spaces. No claim of physical identity, shared dynamics,
universality classes, or scaling behavior is made. The analogy is limited
exclusively to the existence of constraint-defined boundaries across which
classification changes discontinuously.

In physics, phase transitions are characterized by non-analytic behavior induced
by control parameters and thresholds \cite{stanley1971phase}, as well as by
qualitative regime changes in nonlinear systems \cite{strogatz2018nonlinear}.
This paper does not import thermodynamic variables, order parameters, dynamical
laws, or universality arguments from that literature. Instead, it isolates a
single minimal correspondence: the presence of Boolean or threshold predicates
that induce discontinuous changes in admissible classifications.

Within socio-technical and organizational research, related notions such as
resilience, regime shifts, or collapse have been discussed primarily in
descriptive or interpretive terms (e.g., \cite{holling1973resilience}), often
entangled with normative, managerial, or policy-oriented interpretations.
In contrast, the present work remains strictly formal and diagnostic: phases
are defined exclusively as predicate-induced subsets of $\Omega(K_{EA})$, without
prescribing responses, interventions, or adaptive strategies.

The Ontology of Continua (OC Core) and the executable specification
\artifact{ETS.K_EA.v1.1} are treated as internal formal sources
\cite{oc_core_internal,ets_keav11_internal}. They provide the complete and closed
symbol set, predicates, and invariants used in this paper. No external theoretical
framework, empirical model, or sociological theory is imported.

The novelty of the contribution lies solely in the explicit formalization of
phase logic in socio-technical enterprises as a threshold-induced partition of an
admissible state space, together with precise and structurally defined
falsifiability conditions. No claims are made beyond this limited structural
correspondence.
